\section{Temporal and Provenance Workloads (C5--C7)}
\label{sec:temporal_provenance}

\paragraph{Temporal interval queries.}
Event data, financial tickers, and knowledge bases with valid-time semantics require interval-constrained retrieval: given a fact $\phi$ with validity window $[t_s, t_e)$, return all facts whose windows intersect a query interval. In RDF stores, temporal annotations are reified as additional triples, forcing queries to filter on a path expression of length $O(a)$ per fact. SheafDB natively indexes each tuple's temporal attributes in an interval tree backed by augmented B-trees, supporting overlap, contains, and before/after predicates in $O(\log n + k)$ time where $k$ is the result cardinality.

\paragraph{Sliding window aggregation.}
Aggregate computation over sliding windows---e.g., ``sum of sensor readings over the past hour, updated every minute''---requires repeated range queries as the window shifts. A KG approach issues $N$ separate SPARQL queries for $N$ window positions, each performing the full reification traversal. SheafDB executes a single pass over the interval index, maintaining a running aggregate across window positions in $O(n \log n)$ total time, compared to $O(N \cdot n)$ for the SPARQL approach.

\paragraph{Provenance chains and derivation DAGs.}
Provenance workloads (C7) trace the derivation history of a fact through a chain or DAG of transformations. In PROV-O---the W3C provenance ontology---each derivation step is encoded as a reified triple `$\phi \xrightarrow{\text{wasGeneratedBy}} \text{Activity}$' with additional attribution triples. Chaining $d$ steps requires $O(d^2)$ recursive property-path evaluation on most SPARQL engines because each intermediate result must be rejoined with the provenance graph.

SheafDB models provenance as a provenance sheaf: a functor from the derivation poset to the category of tuple sets. Each derivation step is a single lightweight edge in the sheaf rather than a collection of reified triples. Walking a chain of length $d$ costs $O(d)$---linear in the derivation depth---because the sheaf structure preserves the derivation topology without join overhead. Fork--join DAGs are handled by sheaf morphism composition in $O(d + b)$ where $b$ is the branching factor.

\paragraph{Experimental setup.}
Workloads C5 (temporal point-query), C6 (sliding-window aggregation), and C7 (provenance chain traversal) are evaluated at scale 1,000 and 10,000 facts. Complexity analysis for larger scales appears in Section~\ref{sec:complexity}.
\endinput
